\documentclass[varwidth, border = 3pt]{standalone}

\usepackage{main}

%%%
% Nous avons besoin des macros ''\num'' et ''\num''unit \pack
% ''siunitx'', mais comme le \tab est construit par le \pack
% ''tabularray'', nous devons faire appel à la bibliothèque
% ''siunitx'' proposée par ''tabularray''.
%%%
\UseTblrLibrary{siunitx}

\begin{document}

\csvread{mes-donnees}{data.csv}

%%%
% L'option est passée telle quelle à l'\env ''tblr'' de
% ''tabularray''. Nous utilisons ici les possibilités de \mefs
% fines offertes par ''tabularray''.
%
%
% note::
%     ''2-Y'' indique les positions de la \2e à l'avant-dernière,
%     tandis que ''1-Z'' va de la \1ere à la dernière.
%%%
\csvtblr{mes-donnees}[
% Lignes verticales.
  vline{2-Y}   = {solid},
% Lignes horizontales.
  hline{2-Y}   = {solid},
% Cellules de la \1ere ligne contenant du texte en gras.
  cell{1}{1-Z} = {cmd = \bfseries},
% Données de la \2e \col en mode \math.
  cell{2-Z}{2} = {mode = math},
% Données de la \3e \col mises en forme par la macro''\num'' du
% \pack ''siunitx''.
  cell{2-Z}{3} = {cmd = {\num[locale = FR]}},
% Données de la \4e \col mises en forme par la macro''\unit'' du
% \pack ''siunitx''.
  cell{2-Z}{4} = {cmd = \unit}
]

\end{document}
